\documentclass[11pt,a4paper]{article}
\usepackage[utf8]{inputenc}
\usepackage[T1]{fontenc}
\usepackage[english]{babel}
\usepackage{amsmath,amssymb,amsthm}
\usepackage{graphicx}
\usepackage{booktabs}
\usepackage{hyperref}
\usepackage{geometry}
\geometry{margin=2.5cm}

\title{\textbf{SOCRATES Phase 2, Stage 2:} \\
\large Poly-Algebraic Calculus --- Hypergraph Dimension Estimation, \\
a Ten-Problem Physics Benchmark, and Comparison to Classical Methods}
\author{Xavier Callens \\ Socrate AI Lab}
\date{\today}

\begin{document}
\maketitle

\begin{abstract}
This report documents Phase 2 (Stage 2) of the Socrate AI Lab programme:
the construction and rigorous benchmarking of a computational
``Poly-Algebraic Calculus'' toolkit for hypergraph-based physics, following
Wolfram's proposal that space, time, and known physics may emerge from
discrete rewriting of a pre-geometric hypergraph. Of four brainstormed
formalization concepts, two were selected for implementation on tractability
grounds: a higher-arity rule-search calculus (\emph{Poly-Algebraic
Calculus} proper) and an emergent-dimension estimator (\emph{Pre-Geometric
Dimension Variables}). The dimension estimator was applied to ten known
physics problems with independently verifiable dimension targets --- exact
results for closed orbits and rigorous theorems, and literature values
(Grassberger--Procaccia, 1983) for two classical chaotic attractors. The
initial benchmark run produced a headline of ``10/10 passed'' that a
subsequent audit showed to be substantially misleading: six of the ten
passes came from a mathematically degenerate regime, and neither chaotic
result was converged in point count. Two independent, real software defects
were found in the process and are corrected here, with regression tests and
measured before/after evidence, including a $\sim$27$\times$ performance
improvement. A genuine classical baseline (the Grassberger--Procaccia
correlation-sum method) and an apples-to-apples comparison harness were then
built so that ``better than the traditional approach'' could be a checkable
claim rather than an assertion. A systematic re-comparison across all ten
problems, targeting a majority win rate against this baseline, was in
progress at the time of writing and is reported as a preliminary status,
not a final result --- consistent with this programme's standing rule that
an unconverged or incomplete measurement must be labelled as such, never
smoothed into a premature conclusion.
\end{abstract}

\tableofcontents

\section{Introduction and Motivation}

Stephen Wolfram's Physics Project proposes that space, time, and the known
laws of physics are not fundamental, but emerge from the discrete,
repeated application of local rewriting rules to a pre-geometric
hypergraph. Formalizing this hypothesis computationally requires new
mathematical and software tooling, analogous to how historical paradigm
shifts in physics required new mathematical syntax (algebra for the
unknown quantity, calculus for instantaneous change, complex numbers for
solutions ``off the real line'').

Four concepts were brainstormed as candidate formal tools for this
programme:

\begin{enumerate}
\item \textbf{Poly-Algebraic Calculus} (higher-arity algebra): extending
  algebra's ``solve for an unknown quantity'' to ``solve for an unknown
  higher-arity relation or rewrite rule.''
\item \textbf{Branchial Cohomology}: a topological invariant of the
  multiway (branching) graph of all possible rewrite histories, proposed
  as the discrete analogue of quantum interference.
\item \textbf{Rulial Operators}: discrete derivatives measuring the rate
  of topological deformation per rewrite step, proposed as a discrete
  analogue of energy/mass density.
\item \textbf{Pre-Geometric Dimension Variables}: treating the dimension
  of emergent space as a fluid, locally-varying quantity inferred from how
  the hypergraph's ball volume scales with radius.
\end{enumerate}

\subsection{Selection rationale}

Concepts \#1 and \#4 were selected for implementation. Concept \#4
(dimension estimation via volume-growth scaling) is a well-established
technique already used in causal set theory and in Wolfram's own work ---
implementable immediately with no new mathematics required. Concept \#1
(a rule-search calculus) is the substrate every other concept implicitly
depends on: one cannot measure a dimension, a rulial derivative, or a
branchial invariant without first having a rigorous language for
hypergraphs, patterns, and rewrites. Concepts \#2 and \#3 were assessed as
the least tractable: \emph{Branchial Cohomology} in particular has no
working mathematical definition yet beyond an evocative name --- no chain
complex, no boundary operator has been specified for it anywhere in the
literature this programme is aware of, and implementing it would have
required inventing the mathematics from scratch rather than assembling
existing tools, in tension with this programme's discipline of not
asserting more than a computation actually proves.

\section{Methods}

\subsection{Hypergraph representation and rewriting}

A hypergraph is represented as an immutable multiset of ordered hyperedges
over an implicit node set (nodes exist only by virtue of participating in
a relation, following the Wolfram Model convention). Rewrite rules are
specified as a left-hand-side pattern (a tuple of pattern edges over
pattern variables) and a right-hand-side replacement; matching is
performed by backtracking search over consistent pattern-variable
assignments (subhypergraph pattern matching). Two evolution modes are
implemented: a single-history mode that applies all non-overlapping
matches per generation (the standard Wolfram Model update rule), and a
multiway mode that branches over every individual single-match update,
producing the full branching-state graph of the abstract rewriting system.

\subsection{Emergent dimension estimation}

For a hypergraph flattened into its adjacency graph (each hyperedge
becomes a clique on its constituent nodes), the ball of radius $r$ around
a node is the set of nodes reachable within $r$ hops. If the graph behaves
like a $d$-dimensional lattice at some scale, ball volume grows as
$|B(r)| \sim r^d$. The estimator fits \emph{shell size}
($dV(r) = |B(r)| - |B(r-1)|$), not cumulative volume, against radius in
log-log space: cumulative volume on a lattice is affine rather than
homogeneous (a one-dimensional path has $|B(r)| = 2r+1$ from an interior
point, not $r^1$), which biases a naive fit on volume itself well below
the true dimension at any computationally tractable radius. Shell size has
no such additive offset (a path's shell is the exact constant 2; a 2D
grid's shell is exactly $4r$), so the corrected estimator returns
\emph{exactly} 1.0 and 2.0 on these known lattices, not merely
approximately.

\subsection{From continuous trajectories to hypergraphs}

A continuous physical trajectory (a set of points in phase space) is
converted to a hypergraph via a $k$-nearest-neighbour proximity graph: each
point becomes a node, and an edge connects each point to its $k$ nearest
neighbours (symmetrized). Graph distance on a sufficiently dense $k$-NN
graph approximates geodesic distance on the underlying manifold, the
standard justification for this class of dimension estimator
(Grassberger \& Procaccia, 1983; the same construction underlies
Isomap-style manifold learning).

\subsection{The traditional baseline}

To make ``better than the traditional approach'' a checkable claim, the
classical Grassberger--Procaccia correlation integral was implemented as
an explicit reference method:
\[
C(r) = \frac{2}{N(N-1)} \#\{(i,j) : i<j,\ |x_i - x_j| < r\}, \qquad
D_2 = \frac{d \log C(r)}{d \log r}
\]
Both this baseline and the shell-growth estimator are built on the same
underlying primitive (\texttt{scipy.spatial.cKDTree}), so that any
compute-cost comparison isolates the \emph{algorithmic} difference (global
pairwise scaling versus local graph growth) rather than an
implementation-quality difference between a carefully optimized new method
and a deliberately weak baseline.

\subsection{Comparison methodology}

For a given point cloud and known target dimension, each method's minimum
point count for \emph{stable} convergence is determined: the smallest
$n$ in a tested grid at which the estimate is within tolerance of the
target \emph{and remains} within tolerance at every larger $n$ tested ---
guarding against a single lucky crossing of a still-moving estimate, the
exact failure mode discovered in the initial chaotic-attractor results
(Section~\ref{sec:findings}). A ``win'' requires both methods to
demonstrably converge; a method that never converges cannot be credited
with needing ``fewer'' points, since that would reward failure to
converge as though it were efficiency.

\section{Initial Benchmark: Ten Known Physics Problems}
\label{sec:findings}

The dimension estimator was applied to ten problems spanning exact
analytic targets, a rigorous probabilistic theorem, and two literature
values from classical chaos theory.

\begin{table}[h]
\centering
\small
\begin{tabular}{@{}llll@{}}
\toprule
\# & Problem & Known dimension & Source \\
\midrule
1 & Harmonic oscillator & 1 (exact) & closed phase-space curve \\
2 & Nonlinear pendulum (large amplitude) & 1 (exact) & closed periodic orbit \\
3 & Kepler two-body orbit ($e=0.6$) & 1 (exact) & closed ellipse \\
4 & Mars vs.\ real JPL Horizons data & 1 (exact) & smooth arc, real ephemeris \\
5 & Quasi-periodic 2-torus (Lissajous) & 2 (exact) & standard KAM-type result \\
6 & Planar Brownian motion & 2 (exact) & Taylor (1953) theorem \\
7 & Restricted 3-body periodic orbit & 1 (exact, if closure verified) & limit-cycle definition \\
8 & Lorenz attractor & $\approx 2.05$ & Grassberger \& Procaccia (1983) \\
9 & R\"ossler attractor & $\approx 2.01$ & Grassberger \& Procaccia (1983) \\
10 & Driven damped pendulum (mode-locked) & 1 (exact, if non-chaotic) & limit-cycle definition \\
\bottomrule
\end{tabular}
\caption{The ten benchmark problems and their independently-sourced target dimensions.}
\end{table}

\subsection{Headline result and why it is misleading}

All ten problems passed their stated tolerance. An audit pass, followed by
an independent verification of that audit (Section~\ref{sec:audit-chain}),
established that this headline substantially overstates what was actually
demonstrated:

\begin{itemize}
\item \textbf{Six of ten ``passes'' are structurally degenerate.} A $k$-NN
  graph built on any smooth closed curve is an exact circulant ring
  lattice, with a perfectly constant shell size at every radius. The
  estimator is mathematically incapable of returning anything other than
  dimension 1 with $r^2=1.0$ on such an input --- the reported perfect fit
  is a sentinel meaning ``the input had no variation to fit,'' not
  evidence that the method measured anything accurately. This affected
  problems 1, 2, 3, 4, 7, and 10.
\item \textbf{Neither chaotic-attractor result was converged in point
  count.} A convergence sweep (varying the number of sampled points, all
  else held fixed) showed both the Lorenz and R\"ossler dimension
  estimates rising monotonically \emph{through} their literature values
  and past them, still climbing at the largest point count tested. The
  reported ``passing'' errors (0.064 for Lorenz, 0.014 for R\"ossler)
  described where a still-moving quantity happened to be sampled, not a
  converged measurement.
\item \textbf{The two remaining results (Lissajous torus, Brownian
  motion)} were the least-degenerate in the set, but a later correction
  found they had not been checked with the same convergence/sampling-grid
  discipline applied to the chaotic cases, and both showed comparable
  sensitivity to sampling-grid choice once checked.
\end{itemize}

\subsection{Multi-level review chain}
\label{sec:audit-chain}

Ten independent agents each solved one problem and self-reported a result;
all ten self-reports were later confirmed numerically correct. An
independent audit reviewed all ten and correctly identified the structural
problems above. A \emph{second}, independent review was then applied
specifically to the audit's own written analysis --- and found three
genuine errors in it: an over-generalized claim about the scope of a
software defect (Section~\ref{sec:defects}), a speculative reattribution
that incorrectly overturned a separately-correct finding, and a factually
incorrect claim about which sub-results had received a particular check,
which had propped up one of the audit's two headline ``success'' verdicts
without the same scrutiny applied elsewhere. All three errors are
documented and corrected in the underlying benchmark record. This
multi-level structure --- workers, auditor, and an independent check on
the auditor itself --- is now standing practice for any result in this
programme intended to be cited.

\section{Software Defects Found and Corrected}
\label{sec:defects}

Two independent, genuine defects in the estimator's own implementation
were found as a direct consequence of running it against real problems
rather than only the synthetic lattices used in its original unit tests.

\subsection{R\texorpdfstring{$^2$}{²} catastrophic cancellation}

For a perfectly constant shell sequence, the total sum of squares in the
$R^2$ computation should be exactly zero. Floating-point rounding in the
intermediate mean computation left it at approximately $10^{-31}$ for
specific fit lengths and shell values instead of exactly zero, causing the
implementation's exact-zero comparison to take the general-case branch and
divide by that residual --- collapsing the reported $R^2$ to approximately
zero by pure numerical noise, even though the underlying dimension
estimate remained correct. This meant the estimator's own documented
default settings returned an undefined result (\texttt{NaN}) for several
ordinary lattices, including the exact configuration four of the ten
benchmark problems used. The fix replaces the exact-zero comparison with a
relative-scale tolerance; verified against every shell value found to
trigger the defect (6, 17, 18 at the affected fit length) plus negative
controls that were already correct, with two regression tests added.

\subsection{Adjacency-reconstruction performance defect}

A separate, unrelated performance defect was found while building the
comparison harness: the hypergraph's adjacency structure was rebuilt from
scratch, at $O(\text{edges})$ cost, on every single distance query, and
the dimension estimator queried it once per radius per sampled node ---
compounding into thousands of redundant reconstructions for a realistic
convergence sweep. A representative 3200-point sweep took 136 seconds
before correction. The fix caches the adjacency structure per (immutable)
hypergraph instance and restructures the dimension estimator to perform a
single incremental breadth-first search collecting all radii's volumes in
one pass, rather than recomputing from the source at every radius. The
identical 3200-point sweep completed in 5.0 seconds after the fix --- a
$\sim$27$\times$ improvement --- with numerically \emph{identical} results
before and after, confirming the fix traded no correctness for speed. This
performance improvement is the specific enabling condition for the
systematic re-comparison reported as in-progress in
Section~\ref{sec:current-status}: the convergence sweeps it requires were
not practically tractable before this fix.

\subsection{Duplicate-point detection}

A third improvement, not a bug fix but a robustness addition: sampling
multiple periods of a closed orbit produces near-exact duplicate point
clusters (separated by $\sim 10^{-9}$ in the cases observed), which
silently corrupted three of the ten original benchmark results with
severity ranging from a degraded fit quality to a complete failure to
converge. The $k$-NN construction now detects duplicate \emph{clusters}
(not merely pairs, since a three-period orbit produces clusters of three)
via a connected-components analysis on the within-threshold pair graph,
calibrated against the point cloud's bounding-box diagonal rather than its
local nearest-neighbour spacing --- an initial attempt using the latter
failed silently, because a duplicated point's nearest neighbour \emph{is}
its own duplicate, corrupting the very statistic meant to detect it.
Verified on the exact reproduction of the original defect (300 points,
each period-tripled at $\sim 10^{-9}$ offset): the corrected implementation
raises an explicit error naming all 300 size-three clusters, and an
automatic-cleaning option recovers exactly the original 300 points and the
correct dimension of exactly 1.0.

\section{Current Status of the Comparison to Traditional Methods}
\label{sec:current-status}

With the defects above corrected and a genuine traditional-method baseline
implemented (Section 2.4), a systematic re-comparison across all ten
benchmark problems was launched, targeting the programme's stated goal: a
majority (at least eight of ten) of the problems demonstrating an
advantage over the traditional correlation-sum method, defined by two
explicit, checkable criteria rather than a single vague notion of
``better'':

\begin{enumerate}
\item \textbf{Compute savings}: more than 10\% fewer sample points needed
  for equivalent, \emph{stably converged} accuracy against the same
  tolerance.
\item \textbf{Density-robustness} (the concrete, checkable form of
  ``singularity avoidance''): physical trajectories with a close approach
  or a slow/fast turning point (a near-primary passage in the restricted
  three-body problem, perihelion passage in an eccentric Kepler orbit, a
  pendulum's turning point) are sampled far more densely in time near that
  feature than elsewhere, even though the underlying curve is uniformly
  one-dimensional throughout. A \emph{global} method that pools all
  pairwise distances into one statistic (the traditional correlation sum)
  can be measurably biased by such a density cluster; a \emph{local}
  method that only ever examines each point's nearest neighbours (the
  shell-growth estimator) is not, by construction. This is tested
  directly by comparing each method's error against the known target on
  the natural, non-uniformly-sampled data.
\end{enumerate}

\textbf{This comparison was in progress at the time of writing.} One
illustrative, manually-verified data point is available: on a filled
two-dimensional square (target dimension 2, tolerance 0.15), the
traditional correlation-sum method reached stable convergence at 100
points while the shell-growth method required 400, a result favouring the
traditional method on this particular case ($-300\%$ relative point
count, i.e.\ the opposite of a poly-algebraic win). This single spot check
is reported for transparency, not as evidence of the systematic outcome
across all ten problems --- consistent with this programme's standing
rule that a claim resting on a single configuration, rather than a swept
and converged comparison, is not yet a result (Section 2.5). The full,
audited, and independently-verified ten-problem comparison will be
reported in a subsequent revision of this document once complete.

\section{Limitations}

\begin{itemize}
\item \textbf{The estimator has no freedom to be inaccurate in the
  degenerate regime.} Any smooth, densely and evenly sampled closed curve
  will read as dimension 1 with a perfect (but uninformative) fit quality,
  regardless of the method's actual accuracy. A \texttt{degenerate} flag
  now distinguishes this case explicitly in the software, but any result
  drawn from problems 1, 2, 3, 4, 7, or 10 of the original benchmark
  should be read as confirming the pipeline is correctly wired, not as
  evidence of measurement accuracy.
\item \textbf{Fractal (non-integer) dimension targets require a
  convergence sweep, not a single measurement}, and even the corrected
  pipeline's convergence behaviour on the two chaotic attractors in this
  benchmark has not yet been re-established after the defects above were
  fixed --- this is part of the in-progress work of Section~\ref{sec:current-status}.
\item \textbf{The $k$-nearest-neighbour construction remains an
  approximation}: graph distance approximates but does not equal geodesic
  distance on the underlying manifold, and this approximation degrades
  near a manifold's boundary (documented and pinned by a regression test
  showing measurable underestimation at the corner of a sampled cube) and
  is sensitive to the choice of $k$ and to local sample density.
\item \textbf{This work does not validate Wolfram's physical hypothesis.}
  Everything reported here is a statement about the behaviour of a
  specific graph-theoretic algorithm on specific finite point clouds ---
  Tier B by this programme's epistemic convention. No claim is made or
  implied about whether hypergraph rewriting is a correct fundamental
  description of physical spacetime; that remains Tier C (labelled
  conjecture, never load-bearing) throughout.
\end{itemize}

\section{Discussion}

The most consequential finding of this stage is procedural rather than
purely technical: a benchmark against real, independently-sourced targets
found real defects (Section~\ref{sec:defects}) that the estimator's own
unit tests, exercised only against synthetic lattices designed to be easy,
could not have found. Equally consequential is that the \emph{review of
the benchmark} itself required its own independent check before its
conclusions could be trusted --- an audit that correctly identified real
problems in the original work was not, on that basis alone, exempt from
containing real problems of its own. Both observations are recorded as
general lessons for this programme (see the accompanying lessons-learned
document) and are now built into the standing workflow pattern used for
any result in this repository intended to be cited: measurement, audit,
and an independent check on the audit, in that order, with none of the
three treated as sufficient on its own.

\section{Conclusion and Next Steps}

Phase 2, Stage 2 of this programme has produced: (i) a working, tested
implementation of two of the four brainstormed hypergraph-physics
formalization concepts; (ii) a ten-problem benchmark against
independently-sourced physics targets, whose headline result required
substantial correction under audit but whose underlying software defects
were genuinely found and genuinely fixed; (iii) a real classical baseline
and comparison methodology, replacing what would otherwise be an
unfalsifiable claim of superiority with two specific, checkable criteria.
The systematic comparison against that baseline, across all ten problems,
is the immediate next step and was in progress at the time of writing. Its
completion, independently audited and verified to the same standard as
every other result in this document, will be reported as an addendum or
revision rather than assumed in advance.

\section*{References}

\begin{itemize}
\item[] Wolfram, S. (2020). \emph{A Project to Find the Fundamental Theory of Physics.}
\item[] Grassberger, P. \& Procaccia, I. (1983). Measuring the strangeness of
  strange attractors. \emph{Physica D}, 9, 189--208.
\item[] Taylor, S. J. (1953). The Hausdorff $\alpha$-dimensional measure of
  Brownian paths in $n$-space. \emph{Proc. Camb. Phil. Soc.}, 49.
\item[] R\"ossler, O. (1976). An equation for continuous chaos.
  \emph{Phys. Lett. A}, 57.
\item[] Internal references: \texttt{docs/HYPERGRAPH\_NOTES.md},
  \texttt{docs/POLY\_ALGEBRAIC\_BENCHMARK.md}, \texttt{docs/LL.md},
  \texttt{src/socrates/hypergraph/}.
\end{itemize}

\end{document}
